On dispose d'une solution d'acide benzoïque de concentration
$c=\qty{2.5e-2}{\mol\per\L}$.

\begin{enumerate}
  \item L'acide benzoïque, de formule \ch{C6H5COOH}, est un acide faible dans
        l'eau.

        Écrire l'équation de la réaction de l'acide benzoïque avec l'eau.
  \item Calculer le taux d'avancement final de la réaction précédente sachant
        que le pH de la solution est égal à 2,9.
  \item Calculer la conductivité de la solution d'acide benzoïque étudiée.
        Préciser l'unité.
\end{enumerate}
\begin{donnees}
\begin{itemize}
    \item Conductivités molaires ioniques~: \par
    $\lambda_{\ce{H3O+}}=\qty{3,5e-2}{\siemens\square\meter\per\mole}$~;
    $\lambda_{\ce{C6H5COO-}}=\qty{3,2e-3}{\siemens\square\meter\per\mole}$
\end{itemize}
\end{donnees}


